% ── 08 / Why CTO on Demand ──

\section{Why CTO on Demand}

bunq has three options. This section addresses all three directly.

\subsection{Option 1: Build In-House}

Hiring US payment network engineers with experience on a custom core is exceptionally rare. Most US payment engineers have worked exclusively with known cores\dash FIS, Temenos, Thought Machine. Finding the right people while a conditional approval window is running is a multi-month process with no guaranteed outcome.

The deeper issue is the commitment: 3--4 senior engineers at \$180,000--\$250,000 base salary each. \$600,000--\$1,000,000 in annual headcount for a finite integration project. Once the integrations are built, what do those engineers do? A permanent headcount commitment for a time-bounded problem runs against everything bunq stands for. And an engineering hire does not come with compliance expertise\dash that is a separate problem on the same ticking clock.

\subsection{Option 2: Pure Consultancy}

A consultancy advises. bunq still does all the building. Their playbooks and integration experience are based on FIS, Temenos, or similar\dash not a custom core. There is also no sandbox access. A consultancy can tell bunq how to apply for credentials; that is not the same as having the relationships to obtain them.

\subsection{Option 3: CTO on Demand}

\begin{keyinsight}{The Combination That Matters}
Technical execution and regulatory knowledge operate as a single practice. CTO on Demand has built its own custom core banking system. Direct payment network integrations on custom cores are not new territory.
\end{keyinsight}

\begin{itemize}
\item \textbf{Technical}\dash 70--80\% of the work is built against official specifications before sandbox access is required. No dependency on any core vendor.
\item \textbf{Regulatory}\dash Active relationships with the Federal Reserve, OCC, FinCEN, TCH, and NACHA. Sandbox access progressed from week one.
\item \textbf{Compliance}\dash 50+ US-specific templates customised to bunq's model. Independent audit included at no charge.
\item \textbf{Knowledge transfer}\dash CTO on Demand builds alongside bunq's team, trains bunq's engineers throughout, and hands over full documentation and runbooks. bunq owns everything. No ongoing dependency.
\end{itemize}

\subsection{References}

CTO on Demand works with larger fintechs and bank charter applicants across MTL programmes and payment network integrations. References available on request. Morgan and Royce have built these exact integrations in their previous roles at J.P.\ Morgan Chase, American Express, Western Alliance Bank, and USAA.
